% 2007 年山东卷文科第 10 题：程序框图。
%   开始 → 输入 n → S=0,T=0 → 判断 n<2 ?
%   否 → S=S+n → n=n-1 → T=T+n → n=n-1 → 回到判断；
%   是 → 输出 S,T → 结束。
% 来源 PDF：gaokaomath-src/普通高考/2007/2007山东文.pdf
% SHA-256：64b760b73b0a8988bc87dd5d37fc646a10e257f357faf9ad7b1a17fdebaf4af3
% 原页：2007山东文.pdf 第 1 页（A3 三栏）第 2 栏，书页码 399。
% 重绘依据：按原页流程与分支方向重绘，判断框菱形，输入/输出用平行四边形。
\begin{tikzpicture}[scale=1.0, semithick,
  every node/.style={inner sep=2pt},
  node distance=0.55cm,
  box/.style={draw, rounded corners=2pt, minimum width=2.1cm,
    minimum height=0.55cm, align=center},
  para/.style={draw, trapezium, trapezium left angle=70,
    trapezium right angle=110, minimum width=2.3cm, minimum height=0.55cm,
    align=center},
  dec/.style={draw, diamond, aspect=2.4, minimum width=2.6cm,
    minimum height=1.0cm, align=center},
  arr/.style={->, >=stealth}]

  \node[box] (start) at (0,7.4) {开始};
  \node[para] (in) at (0,6.55) {输入 $n$};
  \node[box] (init) at (0,5.70) {$S=0,\ T=0$};
  \node[dec] (test) at (0,4.62) {$n<2?$};
  \node[box] (s) at (0,3.55) {$S=S+n$};
  \node[box] (n1) at (0,2.80) {$n=n-1$};
  \node[box] (t) at (0,2.05) {$T=T+n$};
  \node[box] (n2) at (0,1.30) {$n=n-1$};
  \node[para] (out) at (4.6,4.62) {输出 $S,\ T$};
  \node[box] (end) at (4.6,3.80) {结束};

  \draw[arr] (start) -- (in);
  \draw[arr] (in) -- (init);
  \draw[arr] (init) -- (test);
  \draw[arr] (test) -- node[left] {否} (s);
  \draw[arr] (s) -- (n1);
  \draw[arr] (n1) -- (t);
  \draw[arr] (t) -- (n2);
  % 回到判断
  \draw[arr] (n2.west) -- ++(-1.15,0) |- (test.west);
  % 是分支
  \draw[arr] (test.east) -- node[above] {是} (out.west);
  \draw[arr] (out) -- (end);
\end{tikzpicture}
